% ============================================================
%  示例：智配系统分层架构图
%  基于 harryopo-tikz-diagram layered-arch v2 模板
%  Nature 无衬线风格
%
%  编译命令：xelatex -interaction=nonstopmode example-layered-arch.tex
% ============================================================

\documentclass[12pt,a4paper]{ctexart}

\usepackage{geometry}
\geometry{left=2cm,right=2cm,top=2cm,bottom=2cm}

\usepackage{fontspec}
\usepackage{tikz}
\usetikzlibrary{positioning, fit, backgrounds, arrows.meta, calc}

% ============================================================
% Nature 风格配色
% ============================================================
\definecolor{PaperBorder}{HTML}{B0C4DE}
\definecolor{PaperFill}{HTML}{FFFFFF}
\definecolor{PaperTitle}{HTML}{2C3E50}
\definecolor{PaperMuted}{HTML}{7F8C8D}
\definecolor{AccentOrange}{HTML}{D35400}
\definecolor{PaperLayerBg}{HTML}{F5F8FA}
\definecolor{PaperLayerBorder}{HTML}{D8E3EE}

% ============================================================
% 字体设置 — Nature 无衬线风格
% ============================================================
\setmainfont{Arial}
\setsansfont{Arial}
\setmonofont{Consolas}[Scale=0.90]
\setCJKsansfont{Microsoft YaHei}
\setCJKmainfont{Microsoft YaHei}

% ============================================================
% 间距常量
% ============================================================
\def\ModW{4.0cm}
\def\ModH{1.5cm}
\def\WideH{0.95cm}
\def\ModGap{0.55cm}
\def\InterGap{1.6cm}
\def\IntraGap{0.7cm}
\def\LayerPad{12pt}

\begin{document}

\begin{center}
  {\Large\sffamily\bfseries\textcolor{PaperTitle}{智配系统分层架构图}}\\[0.3em]
  {\small\sffamily\itshape\textcolor{PaperMuted}{Zhipei System --- Layered Architecture}}
\end{center}

\vspace{1em}

\begin{figure}[htbp]
\centering
\begin{tikzpicture}[
    node distance = 0pt and \ModGap,
    >=Stealth,
    line width = 0.7pt,
    every node/.style = {inner sep=0pt, outer sep=0pt},
    % ---- 模块节点 ----
    module/.style = {
        rectangle,
        draw = PaperBorder,
        fill = PaperFill,
        line width = 0.8pt,
        text width = \ModW,
        minimum height = \ModH,
        font = \small\sffamily,
        align = center,
        rounded corners = 5pt,
        inner sep = 8pt,
        text = PaperTitle,
    },
    module-wide/.style = {
        module,
        text width = 3*\ModW + 2*\ModGap,
        minimum height = \WideH,
    },
    mod-title/.style = {
        font = \small\sffamily\bfseries,
        text = PaperTitle,
    },
    mod-desc/.style = {
        font = \footnotesize\sffamily,
        text = PaperMuted,
    },
    mod-tech/.style = {
        font = \footnotesize\sffamily\ttfamily,
        text = AccentOrange,
    },
    % ---- 层容器 ----
    layer-box/.style = {
        draw = PaperLayerBorder,
        fill = PaperLayerBg,
        line width = 0.7pt,
        dashed,
        dash pattern = on 4pt off 3pt,
        rounded corners = 8pt,
        inner sep = \LayerPad,
    },
    layer-title/.style = {
        font = \normalsize\sffamily\bfseries,
        text = PaperTitle,
        anchor = north west,
        inner sep = 0pt,
    },
    layer-tag/.style = {
        font = \footnotesize\sffamily\itshape,
        text = PaperMuted,
        anchor = north east,
        inner sep = 0pt,
    },
    % ---- 箭头 ----
    arr-down/.style = {
        ->, >=Stealth, thick, PaperBorder,
    },
    arr-dep/.style = {
        ->, >=Stealth, thick, PaperMuted,
    },
    arr-label/.style = {
        font = \footnotesize\sffamily\bfseries,
        fill = white,
        text = PaperBorder,
        inner sep = 2pt,
        outer sep = 0pt,
    },
    % ---- 图例 ----
    legend-box/.style = {
        draw = PaperLayerBorder,
        fill = PaperFill,
        line width = 0.6pt,
        rounded corners = 5pt,
        inner sep = 8pt,
        font = \footnotesize\sffamily,
        text = PaperTitle,
    },
]

    % ============================================================
    %  第一层：智配桌面安装器（含宽模块）
    % ============================================================

    \node[module] (L1a) {
        {\mod-title 硬件检测模块}\\[2pt]
        {\mod-tech GPU / VRAM / RAM}\\[1pt]
        {\mod-tech CPU / OS}
    };

    \node[module, right=\ModGap of L1a] (L1b) {
        {\mod-title 模型推荐引擎}\\[2pt]
        {\mod-desc 根据硬件推荐}\\[1pt]
        {\mod-desc 最佳模型方案}
    };

    \node[module, right=\ModGap of L1b] (L1c) {
        {\mod-title 一键部署器}\\[2pt]
        {\mod-desc 自动化安装}\\[1pt]
        {\mod-tech Ollama + WebUI}
    };

    \node[module-wide, below=\IntraGap of L1a.south west, anchor=north west] (L1w) {
        {\mod-title 管理后台（Web Dashboard）}\\[2pt]
        {\mod-desc 模型管理\quad$\bullet$\quad 运维监控\quad$\bullet$\quad 知识库\quad$\bullet$\quad 日志}
    };

    \begin{scope}[on background layer]
        \node[layer-box, fit=(L1a)(L1b)(L1c)(L1w)] (box1) {};
    \end{scope}

    \node[layer-title] at ($(box1.north west)+(0,4pt)$) {智配桌面安装器};
    \node[layer-tag] at ($(box1.north east)+(0,4pt)$) {Electron 桌面应用};

    % ============================================================
    %  第二层：后端服务
    % ============================================================

    \node[module, below=\InterGap of L1w.south west, anchor=north west] (L2a) {
        {\mod-title 部署引擎}\\[2pt]
        {\mod-tech Ollama 管理}\\[1pt]
        {\mod-desc 模型生命周期}
    };

    \node[module, right=\ModGap of L2a] (L2b) {
        {\mod-title 运维监控}\\[2pt]
        {\mod-tech GPU / 内存采样}\\[1pt]
        {\mod-desc 进程守护与告警}
    };

    \node[module, right=\ModGap of L2b] (L2c) {
        {\mod-title API 服务}\\[2pt]
        {\mod-tech FastAPI 框架}\\[1pt]
        {\mod-tech RESTful 接口}
    };

    \begin{scope}[on background layer]
        \node[layer-box, fit=(L2a)(L2b)(L2c)] (box2) {};
    \end{scope}

    \node[layer-title] at ($(box2.north west)+(0,4pt)$) {后端服务};
    \node[layer-tag] at ($(box2.north east)+(0,4pt)$) {Python};

    % ============================================================
    %  第三层：开源生态
    % ============================================================

    \node[module, below=\InterGap of L2a.south west, anchor=north west] (L3a) {
        {\mod-title \ttfamily Ollama}\\[2pt]
        {\mod-desc 本地大模型推理引擎}\\[1pt]
        {\mod-desc 多模型管理}
    };

    \node[module, right=\ModGap of L3a] (L3b) {
        {\mod-title Open WebUI}\\[2pt]
        {\mod-desc 可扩展聊天界面}\\[1pt]
        {\mod-tech RAG / 多模态}
    };

    \node[module, right=\ModGap of L3b] (L3c) {
        {\mod-title 模型文件}\\[2pt]
        {\mod-tech Qwen / Llama 等}\\[1pt]
        {\mod-tech GGUF 量化格式}
    };

    \begin{scope}[on background layer]
        \node[layer-box, fit=(L3a)(L3b)(L3c)] (box3) {};
    \end{scope}

    \node[layer-title] at ($(box3.north west)+(0,4pt)$) {开源生态};
    \node[layer-tag] at ($(box3.north east)+(0,4pt)$) {底层依赖};

    % ============================================================
    %  层间箭头
    % ============================================================

    \draw[arr-down] (box1.south) -- (box2.north)
        node[midway, arr-label, anchor=west, xshift=3pt] {调用};

    \draw[arr-dep] (box2.south) -- (box3.north)
        node[midway, arr-label, anchor=west, xshift=3pt] {依赖};

    % ============================================================
    %  图例（嵌入右下角）
    % ============================================================

    \node[legend-box, anchor=north east] at ($(box3.south east)+(0,-0.8cm)$) {
        \begin{tabular}{@{}c@{\hspace{6pt}}l@{\hspace{14pt}}c@{\hspace{6pt}}l@{}}
            \textcolor{PaperBorder}{\rule[1pt]{14pt}{1.4pt}$>$} & {\sffamily\bfseries\small 调用关系} &
            \textcolor{PaperMuted}{\rule[1pt]{14pt}{1.4pt}$>$} & {\sffamily\bfseries\small 底层依赖}
        \end{tabular}
    };

\end{tikzpicture}
\caption{\sffamily 智配系统分层架构图}
\label{fig:zhipei-layered-arch}
\end{figure}

\end{document}
